\chapter{التفاعلات الأساسية}\label{ch:g11:fundamental-interactions}

ادلُك بالونًا على شعرك فيلتقط قصاصات ورق --- متغلبًا،
ببضعة سنتيمترات مربعة من المطاط، على جذب الكوكب كله
الثقالي. وذلك النصر السهل هو سبب تماسك الذرات وسبب تحطّم
النوى الثقيلة في نهاية المطاف (\cref{ch:g11:nucleus-radioactivity}).
ويسير الكون على أربعة تفاعلات بالضبط؛ ويلقاها هذا الفصل،
ويزنها بعضها ببعض، ويعيّن لكلٍّ طابقه.

\section{المادة من الكواركات إلى المجرّات}

\begin{definition}[سُلّم البنية]\label{def:g11:fundamental-interactions:ladder}
تُبنى المادة طوابق، كلٌّ مركَّب من الذي تحته: فثلاثة
\emph{كواركات}\index{كوارك} (نقطية، أصغر من \qty{e-18}{m}) ترتبط
في \emph{نوكليون}\index{نوكليون} --- بروتون أو نيوترون، قطره نحو
\qty{e-15}{m}؛ وتنرصّ النوكليونات في \emph{نواة}\index{نواة}
(\qty{e-15}{m} إلى \qty{e-14}{m})؛ والنواة مع سحابة إلكتروناتها
\emph{ذرة}\index{ذرة} (\qty{e-10}{m})؛ وترتبط الذرات في
\emph{جزيئات}\index{جزيء} وبلورات (\qty{e-9}{m} فما فوق)، وهذه
تتركّب في أجسام يومية (\qty{1}{m})، وكواكب (\qty{e7}{m})،
ومنظومات كوكبية (\qty{e11}{m})، ومجرّات (\qty{e21}{m})، وفي
الكون المرصود (\qty{e26}{m}).
\end{definition}

\begin{definition}[التفاعلات الأربعة]\label{def:g11:fundamental-interactions:four}
\emph{التفاعل الأساسي}\index{تفاعل أساسي} \omterm{def:g10:inertia:force}{قوة}
لا تُردّ إلى قوى أخرى؛ وهي أربعة بالضبط:
\emph{الجاذبية}\index{الجاذبية}، و\emph{التفاعل الكهرمغناطيسي}\index{تفاعل كهرمغناطيسي}،
\omterm{def:g11:fundamental-interactions:strong}{والتفاعل النووي القوي}، \omterm{def:g11:fundamental-interactions:weak}{والتفاعل النووي الضعيف}.
\end{definition}

\begin{remark}\label{rem:g11:fundamental-interactions:empty}
تفصل بين الطوابق فراغات: \omterm{def:g11:fundamental-interactions:ladder}{فالذرة} أوسع $10^{5}$ مرة من
نواتها --- كبّر \omterm{def:g11:fundamental-interactions:ladder}{النواة} إلى بلية قطرها \qty{1}{cm} فتصير
سحابة الإلكترونات كيلومترًا عرضًا. فلا بد أن شيئًا يمسك هذه الطوابق
المتخلخلة، ولا تكاد \omterm{def:g11:fundamental-interactions:four}{الجاذبية} تكونه أبدًا.
\end{remark}

\begin{omfigure}
\begin{tikzpicture}[font=\small, x=0.24cm]
  \draw[-Stealth] (-1,0) -- (46.5,0);
  \foreach \x/\l/\d in {0/{10^{-18}}/2pt, 3/{10^{-15}}/13pt,
                        8/{10^{-10}}/2pt, 18/{10^{0}}/2pt,
                        25/{10^{7}}/2pt, 29/{10^{11}}/13pt,
                        39/{10^{21}}/2pt, 44/{10^{26}}/13pt}
    \draw (\x,0.1) -- (\x,-0.1) node[below=\d] {$\l$};
  \foreach \x/\n in {0/quark, 3/nucleus, 8/atom, 18/human, 25/Earth,
                     29/{Sun--Earth}, 39/galaxy, 44/universe}{
    \fill[omDef] (\x,0) circle (1.6pt);
    \node[rotate=45, anchor=west, inner sep=1.5pt] at (\x,0.24) {\n};}
\end{tikzpicture}

{\small سُلّم البنية على محور بقوى العشرة: أربعة وأربعون
عقدًا تفصل \omterm{def:g11:fundamental-interactions:ladder}{الكوارك} عن الكون المرصود. والأحجام
\omterm{def:g10:orders-of-magnitude:unit}{بالمتر}.}
\end{omfigure}

\section{الشحنة الكهربائية}

\begin{definition}[الشحنة الكهربائية]\label{def:g11:fundamental-interactions:charge}
\emph{الشحنة الكهربائية}\index{شحنة كهربائية} هي خاصية المادة
التي يؤثر فيها \omterm{def:g11:fundamental-interactions:four}{التفاعل الكهرمغناطيسي}، كما أن الكتلة هي الخاصية
التي تؤثر فيها \omterm{def:g11:fundamental-interactions:four}{الجاذبية}. وتقاس \omterm{def:g10:orders-of-magnitude:unit}{بوحدة} \emph{الكولوم}\index{كولوم}
(\unit{C}) ولها \emph{إشارتان}، موجبة وسالبة، تتلاشيان
عند الجمع؛ والشحنتان المتماثلتان تتنافران، والمختلفتان تتجاذبان.
\end{definition}

\begin{definition}[الشحنة الأولية]\label{def:g11:fundamental-interactions:elementary}
الشحنة محبَّبة: فكل شحنة مرصودة مضاعف صحيح للمقدار المسمّى
\emph{الشحنة الأولية}\index{شحنة أولية}
$e = \qty{1.60e-19}{C}$. ويحمل البروتون $+e$، والإلكترون $-e$، و
النيوترون $0$؛ والشحنة $q$ تحوي $N = \abs{q}/e$ شحنة أولية.
\end{definition}

\begin{proposition}[انحفاظ الشحنة]\label{prop:g11:fundamental-interactions:conservation}
الشحنة الكلية لجملة معزولة لا تتغير أبدًا: فالشحنة لا تُخلق
ولا تُفنى، إنما تُنقل --- وعمليًّا بالإلكترونات،
وهي أخفّ الحاملات.
\end{proposition}
\admitted

\begin{remark}\label{rem:g11:fundamental-interactions:conservation}
لم يُرصد أي خرق قط؛ أما السبب العميق، وهو تناظر في
\omterm{prop:g10:signals-and-waves:em}{الكهرمغناطيسية}، فيُعطى في المجلدات الجامعية.
\end{remark}

\begin{definition}[الشحن بالدلك وبالتماس]\label{def:g11:fundamental-interactions:charging}
دلك عازلين ينزع إلكترونات من أحدهما إلى الآخر،
فيترك شحنتين متعاكستين متساويتي الشدة (\emph{الشحن
بالدلك}\index{الشحن بالدلك})؛ وملامسة جسم مشحون لجسم
متعادل تتقاسم الشحنة (\emph{الشحن
بالتماس}\index{الشحن بالتماس}). ولا تنتقل إلا الإلكترونات؛ والشحنة
الكلية محفوظة.
\end{definition}

\begin{example}[عدّ الإلكترونات]\label{ex:g11:fundamental-interactions:balloon}
بالون مدلوك على شعر يكتسب $q = \qty{-1.6e-8}{C}$: أي أنه
\emph{\omterm{def:g10:signals-and-waves:oscilloscope}{كسب}} $N = \num{1.6e-8} / \num{1.6e-19} = \num{1.0e11}$
إلكترونًا، ويبقى للشعر $+\qty{1.6e-8}{C}$ بالضبط.
\end{example}

\section{قانون كولوم}

\begin{theorem}[قانون كولوم]\label{thm:g11:fundamental-interactions:coulomb}
شحنتان نقطيتان $q_1$ و$q_2$ تفصل بينهما مسافة $d$ تمارس كلٌّ
على الأخرى \omterm{def:g10:inertia:force}{قوة} على المستقيم الواصل بينهما، وشدتاهما متساويتان
واتجاهاهما متعاكسان --- متنافرتان في حالة الإشارتين المتماثلتين، متجاذبتان في حالة الإشارتين
المختلفتين --- حيث
\[
F = k\,\frac{\abs{q_1 q_2}}{d^{2}},
\qquad k = \qty{8.99e9}{N.m^2/C^2}.
\]
\end{theorem}
\admitted

\begin{remark}\label{rem:g11:fundamental-interactions:torsion}
أرسى كولوم القانون سنة 1785 بقياس فتل خيط فتل رفيع؛ أما لماذا يكون \omterm{def:g10:orders-of-magnitude:scinot}{الأس} $2$ بالضبط فيُجاب عنه في المجلدات الجامعية.
\end{remark}

\begin{omfigure}
\begin{tikzpicture}[font=\small]
  \draw[omDef, thick, fill=omDef!20] (0,0) circle (0.3) node[black] {$+$};
  \draw[omDef, thick, fill=omDef!20] (2.6,0) circle (0.3) node[black] {$+$};
  \draw[-Stealth, omThm, very thick] (-0.3,0) -- (-1.5,0)
    node[midway, above] {$\vect F_{2/1}$};
  \draw[-Stealth, omThm, very thick] (2.9,0) -- (4.1,0)
    node[midway, above] {$\vect F_{1/2}$};
\end{tikzpicture}\qquad
\begin{tikzpicture}[font=\small]
  \draw[omDef, thick, fill=omDef!20] (0,0) circle (0.3) node[black] {$+$};
  \draw[omProp, thick, fill=omProp!25] (2.6,0) circle (0.3) node[black] {$-$};
  \draw[-Stealth, omThm, very thick] (0.3,0) -- (1.1,0)
    node[midway, above] {$\vect F_{2/1}$};
  \draw[-Stealth, omThm, very thick] (2.3,0) -- (1.5,0)
    node[midway, above] {$\vect F_{1/2}$};
\end{tikzpicture}

{\small الشحنتان المتماثلتان تتنافران (يمينًا)، والمختلفتان تتجاذبان (يسارًا)؛
وللقوتين دائمًا شدتان متساويتان واتجاهان متعاكسان.}
\end{omfigure}

\begin{remark}[توأم صوري للجاذبية]\label{rem:g11:fundamental-interactions:twin}
قانون كولوم هو، رمزًا برمز، قانون \omterm{def:g11:fundamental-interactions:four}{الجاذبية} الكونية
(\cref{ch:g10:universal-gravitation}) بالشحن مكان الكتل:
\begin{center}
\begin{tabular}{lll}
المنبع & الكتلة $m$ & الشحنة $q$ \\
القانون & $F = G\,\dfrac{m_1 m_2}{d^2}$ & $F = k\,\dfrac{\abs{q_1 q_2}}{d^2}$ \\[1.5ex]
الثابت & $G = \qty{6.67e-11}{N.m^2/kg^2}$ & $k = \qty{8.99e9}{N.m^2/C^2}$ \\
\omterm{def:g10:signals-and-waves:signal}{الإشارة} & جاذبة دائمًا & جاذبة أو نافرة \\
\end{tabular}
\end{center}
والفرقان يقودان كل ما يأتي أدناه: فالكهرباء أقوى بفارق
هائل، وهي وحدها تقدر أن تلغي نفسها.
\end{remark}

\begin{example}[بروتونان والقوتان معًا]\label{ex:g11:fundamental-interactions:protons}
بروتونان ($m_p = \qty{1.67e-27}{kg}$) يقعان على بعد $d = \qty{1.0e-15}{m}$
--- أي جارين نوويين. والتنافر الكهربائي:
\[
F_E = \num{8.99e9} \times
\frac{(\num{1.60e-19})^{2}}{(\num{1.0e-15})^{2}}
\approx \qty{2.3e2}{N}
\]
--- أي \omterm{def:g10:universal-gravitation:weight}{وزن} حقيبة كتلتها $23$ كيلوغرامًا، يحملها جسيم كتلته
$10^{-27}$ \omterm{def:g10:orders-of-magnitude:unit}{كيلوغرام}. أما التجاذب الثقالي:
$F_G = \num{6.67e-11} \times (\num{1.67e-27})^{2} /
(\num{1.0e-15})^{2} \approx \qty{1.9e-34}{N}$. والنسبة
$F_E / F_G = k e^{2} / (G m_p^{2}) \approx \num{1.2e36}$
مستقلة عن المسافة (فتتلاشى $d^{2}$): أي أن \omterm{def:g11:fundamental-interactions:four}{الجاذبية} بين الجسيمات
الأولية أضعف $10^{36}$ مرة من أن تهمّ --- على \emph{أي} مسافة.
\end{example}

\section{التفاعلان النوويان القوي والضعيف}

يترك المثال فضيحة: \omterm{def:g11:fundamental-interactions:ladder}{فنواة} الهيليوم تمسك بروتونين
يتنافران بعشرات النيوتنات --- ومع ذلك يوجد الهيليوم.

\begin{definition}[التفاعل النووي القوي]\label{def:g11:fundamental-interactions:strong}
\emph{التفاعل النووي القوي}\index{تفاعل نووي قوي} يربط \omterm{def:g11:fundamental-interactions:ladder}{الكواركات}
في نوكليونات والنوكليونات في نوى. وعند \qty{e-15}{m} يكون
جاذبًا وأقوى بكثير من التنافر الكهربائي (بآلاف
النيوتنات بين النوكليونات)؛ ووراء بضعة أمتار من رتبة $10^{-15}$ يتلاشى:
إذ إنّ \emph{مداه}\index{مدى (تفاعل)} نحو
\qty{e-15}{m}. وهو يمسك البروتونات والنيوترونات سواءً ويتجاهل الشحنة.
\end{definition}

\begin{omfigure}
\begin{tikzpicture}[font=\small]
  \draw[omDef, thick, fill=omDef!20] (0,0) circle (0.34) node[black] {p};
  \draw[omDef, thick, fill=omDef!20] (3.2,0) circle (0.34) node[black] {p};
  \draw[-Stealth, omThm, thick] (-0.34,0) -- (-1.7,0)
    node[midway, above] {$F_E \approx \qty{230}{N}$};
  \draw[-Stealth, omThm, thick] (3.54,0) -- (4.9,0) node[midway, above] {$F_E$};
  \draw[-Stealth, omMeth, very thick] (0.2,0.55) -- (1.45,0.55)
    node[midway, above] {النووية القوية};
  \draw[-Stealth, omMeth, very thick] (3.0,0.55) -- (1.75,0.55)
    node[midway, above] {النووية القوية};
  \draw[Stealth-Stealth] (0,-0.65) -- (3.2,-0.65)
    node[midway, below] {$d = \qty{1.0e-15}{m}$};
\end{tikzpicture}

{\small شدّ الحبل داخل \omterm{def:g11:fundamental-interactions:ladder}{النواة}: فعند $10^{-15}$~m يغلب التجاذب
القوي (آلاف النيوتنات) التنافرَ الكهربائي.}
\end{omfigure}

\begin{remark}[لماذا توجد النوى --- ولماذا تستسلم الثقيلة منها]\label{rem:g11:fundamental-interactions:nuclei}
تتغير القوتان في \omterm{def:g11:fundamental-interactions:ladder}{النواة} على نحوين مختلفين. فالغراء القوي
قصير المدى: إذ لا يربط \omterm{def:g11:fundamental-interactions:ladder}{النوكليون} إلا جيرانه القلائل ضمن
\qty{e-15}{m}. أما التنافر الكهربائي فطويل المدى: فكل بروتون
يدفع \emph{كل} بروتون آخر. وكلما كبرت النوى تراكم التنافر
أسرع من الغراء: فتحتاج النوى الثقيلة نيوترونات إضافية (غراء
بلا شحنة)، ووراء نحو $80$ بروتونًا يخفق حتى ذلك --- فتعيش
أثقل النوى على وقت مستعار (\cref{ch:g11:nucleus-radioactivity}).
\end{remark}

\begin{definition}[التفاعل النووي الضعيف]\label{def:g11:fundamental-interactions:weak}
\emph{التفاعل النووي الضعيف}\index{تفاعل نووي ضعيف}، \omterm{def:g11:fundamental-interactions:strong}{ومداه} أقصر
بعدُ، يتيح لنيوترون أن يصير بروتونًا (والعكس): فهو يقود
النشاط الإشعاعي $\beta$ في \cref{ch:g11:nucleus-radioactivity}.
\end{definition}

\section{من يحكم أي طابق}

\begin{method}[تدقيق سُلّم]\label{met:g11:fundamental-interactions:audit}
لمعرفة أي تفاعل يحكم بنية حجمها $d$، تحقّق من المدى
ومن الإلغاء:
\begin{enumerate}
  \item $d \leq \qty{e-14}{m}$: يحكم \emph{\omterm{def:g11:fundamental-interactions:strong}{التفاعل النووي القوي}} ---
        فهو يغلب التنافر الكهربائي، \omterm{def:g11:fundamental-interactions:four}{والجاذبية} $10^{36}$ خارج اللعبة.
  \item من الذرات إلى الجبال (من $\qty{e-10}{m}$ إلى $\qty{e4}{m}$):
        \omterm{def:g10:inertia:force}{القوة} النووية القوية خارج المدى؛ فيحكم \emph{\omterm{def:g11:fundamental-interactions:four}{التفاعل
        الكهرمغناطيسي}} --- أي الروابط والتماسك \omterm{def:g10:inertia:inventory}{والاحتكاك} والتماس.
  \item الكواكب فما فوق ($d \geq \qty{e6}{m}$): تكون المادة متعادلة بدقة
        هائلة، فتلغي الكهرباء نفسها؛ أما الكتلة فلها
        \omterm{def:g10:signals-and-waves:signal}{إشارة} واحدة وتتجمع دائمًا --- فتحكم \emph{\omterm{def:g11:fundamental-interactions:four}{الجاذبية}} السماء.
\end{enumerate}
أما التفاعل الضعيف فلا يحكم طابقًا: فهو لا يربط شيئًا، لكنه يحكّم في
التحولات (تفكك $\beta$).
\end{method}

\begin{omfigure}
\begin{tikzpicture}[font=\small, x=0.26cm]
  \draw[-Stealth] (-1,0) -- (41.5,0);
  \foreach \x/\l in {3/{10^{-15}}, 8/{10^{-10}}, 18/{10^{0}},
                     24/{10^{6}}, 39/{10^{21}}}
    \draw (\x,0.1) -- (\x,-0.1) node[below] {$\l$};
  \draw[omProp, fill=omProp!35] (3,0.3) rectangle (4,0.75);
  \node[anchor=west, omProp] at (4.4,0.52) {النووية القوية};
  \draw[omDef, fill=omDef!25] (8,0.95) rectangle (22,1.4);
  \node[omDef] at (15,1.17) {الكهرمغناطيسية};
  \draw[omThm, fill=omThm!25] (24,1.6) rectangle (39,2.05);
  \node[omThm] at (31.5,1.82) {الجاذبية};
\end{tikzpicture}

{\small من يحكم أين: القوي للنوى، والكهرمغناطيسي من الذرات
إلى الجبال، \omterm{def:g11:fundamental-interactions:four}{والجاذبية} للكواكب فما فوق. والأحجام \omterm{def:g10:orders-of-magnitude:unit}{بالمتر}.}
\end{omfigure}

\section{تمارين}

\begin{exercise}[$\star$]\label{exo:g11:fundamental-interactions:1}
أعطِ رتبة قدر كل حجم مما يلي وطابقه من السُّلّم:
البروتون \qty{1.7e-15}{m}؛ \omterm{def:g11:fundamental-interactions:ladder}{وجزيء} الماء \qty{2.8e-10}{m}؛ وحبة الرمل
\qty{2e-4}{m}؛ والأرض \qty{1.3e7}{m}؛ ودرب التبانة \qty{9.5e20}{m}.
\end{exercise}

\begin{exercise}[$\star$]\label{exo:g11:fundamental-interactions:2}
بالون مدلوك على شعر يحمل $q = \qty{-4.8e-9}{C}$. أكسب
إلكترونات أم فقدها، وكم؟ وما شحنة الشعر بعد ذلك،
وأيّ قانون يقول ذلك؟
\end{exercise}

\begin{exercise}[$\star$]\label{exo:g11:fundamental-interactions:3}
شحنتان $q_1 = \qty{+2.0e-6}{C}$ و$q_2 = \qty{-3.0e-6}{C}$ تفصل بينهما
\qty{30}{cm}. احسب \omterm{def:g10:inertia:force}{القوة}؛ أجاذبة هي أم نافرة؟
وقارن \omterm{def:g10:inertia:force}{القوة} على $q_1$ \omterm{def:g10:inertia:force}{بالقوة} على $q_2$.
\end{exercise}

\begin{exercise}[$\star$]\label{exo:g11:fundamental-interactions:4}
تتجاذب شحنتان \omterm{def:g10:inertia:force}{بقوة} $F_0$. فماذا تصير إذا: تضاعفت
المسافة؛ أو قُسمت المسافة على $3$؛ أو تضاعفت إحدى الشحنتين؛
أو تضاعفت الشحنتان \emph{والمسافة} معًا؟
\end{exercise}

\begin{exercise}[$\star$]\label{exo:g11:fundamental-interactions:5}
سمِّ التفاعل المسؤول عن: (أ) دوران القمر حول الأرض؛
(ب) تماسك بلورة ملح؛ (ج) تماسك \omterm{def:g11:fundamental-interactions:ladder}{نواة}
هيليوم؛ (د) تفكك $\beta$ للكربون-14؛ (ه) التصاق بالون
\cref{exo:g11:fundamental-interactions:2} بجدار.
\end{exercise}

\begin{exercise}[$\star\star$]\label{exo:g11:fundamental-interactions:6}
يقع بروتونا \omterm{def:g11:fundamental-interactions:ladder}{نواة} هيليوم على بعد $d = \qty{2.0e-15}{m}$
($m_p = \qty{1.67e-27}{kg}$). احسب تنافرهما الكهربائي، وتجاذبهما
الثقالي، والنسبة. فما الذي يبقي \omterm{def:g11:fundamental-interactions:ladder}{النواة} سليمة؟
\end{exercise}

\begin{exercise}[$\star\star$]\label{exo:g11:fundamental-interactions:7}
في \omterm{def:g11:fundamental-interactions:ladder}{ذرة} هيدروجين، يقع الإلكترون ($m_e = \qty{9.11e-31}{kg}$) على بعد
$d = \qty{5.3e-11}{m}$ من البروتون ($m_p = \qty{1.67e-27}{kg}$).
احسب القوتين الكهربائية والثقالية بينهما ونسبتهما.
فما الذي يمسك الذرات معًا؟
\end{exercise}

\begin{exercise}[$\star\star$]\label{exo:g11:fundamental-interactions:8}
تحمل الكرة A شحنة $q_A = \qty{+8.0e-9}{C}$؛ وكرة B مماثلة لها
متعادلة. تُلامَسان معًا، ثم تُفصلان إلى \qty{10}{cm}.
فأيّ شحنة تحمل كلٌّ منهما (تحقّق من انحفاظ الشحنة)؟ واحسب
\omterm{def:g10:inertia:force}{القوة} بينهما --- أجاذبة أم نافرة؟
\end{exercise}

\begin{exercise}[$\star\star$]\label{exo:g11:fundamental-interactions:9}
شحنتان متماثلتان $q = \qty{1.0e-6}{C}$ تتنافران \omterm{def:g10:inertia:force}{بقوة} \qty{0.90}{N}:
فكم المسافة بينهما؟ وأين تهبط \omterm{def:g10:inertia:force}{القوة} إلى \qty{0.10}{N}؟
\end{exercise}

\begin{exercise}[$\star\star$]\label{exo:g11:fundamental-interactions:10}
بروتونان في جزيئين متجاورين يقعان على بعد $d = \qty{1.0e-10}{m}$.
احسب تنافرهما الكهربائي. وبماذا يسهم \omterm{def:g11:fundamental-interactions:strong}{التفاعل النووي القوي}
على هذه المسافة؟ وأيّ تفاعل يربط \omterm{def:g11:fundamental-interactions:ladder}{الجزيئات}؟
\end{exercise}

\begin{exercise}[$\star\star$]\label{exo:g11:fundamental-interactions:11}
تحوي كل قطعة نقود من قطعتين نحو $\num{e23}$ إلكترون. انقل إلكترونًا
واحدًا من كل \emph{مليون} بينهما وضع القطعتين على بعد \qty{1.0}{m}.
احسب شحنة كل قطعة، ثم \omterm{def:g10:inertia:force}{القوة}. ووزنُ أيّ كتلة
يساويها؟ واستنتج بشأن تعادل المادة اليومية.
\end{exercise}

\begin{exercise}[$\star\star\star$]\label{exo:g11:fundamental-interactions:12}
\omterm{def:g11:fundamental-interactions:ladder}{نواة} يورانيوم قطرها \qty{1.4e-14}{m} وفيها $92$ بروتونًا.
احسب التنافر بين بروتونين على طرفيها المتقابلين. فهل يقدر
\omterm{def:g11:fundamental-interactions:strong}{التفاعل النووي القوي} على ربط هذا الزوج --- ولماذا لا؟ وكم زوجًا
متنافرًا من البروتونات فيها؟ ولماذا تحتاج النوى الثقيلة نيوترونات إضافية،
ولماذا لا تقدر حتى النيوترونات على إنقاذها وراء حجم ما؟
\end{exercise}

\begin{exercise}[$\star\star\star$]\label{exo:g11:fundamental-interactions:13}
كرتان كتلة كلٍّ منهما \qty{1.0}{g} معلّقتان بخيطين طول كلٍّ منهما \qty{50}{cm} مربوطين
بالنقطة نفسها. وإذا أُعطيتا الشحنة $q$ نفسها، استقرتا على بعد \qty{6.0}{cm}
إحداهما من الأخرى. وتكون كل كرة في توازن تحت وزنها \omterm{def:g10:inertia:inventory}{وشدّ} الخيط
\omterm{def:g10:inertia:force}{والقوة} الكهربائية: بيّن أن $F = mg\tan\theta$ (حيث $\theta$:
زاوية الخيط مع العمودي) وأن $\tan\theta \approx 0.060$ هنا.
احسب $F$، ثم $q$، ثم عدد الشحنات الأولية المنقولة.
\end{exercise}

\begin{exercise}[$\star\star\star$]\label{exo:g11:fundamental-interactions:14}
المعطيات: $M_E = \qty{5.97e24}{kg}$، $M_M = \qty{7.35e22}{kg}$، والمسافة بين الأرض
والقمر $d = \qty{3.84e8}{m}$. احسب \omterm{def:g10:universal-gravitation:force}{قوة الجاذبية} بين الأرض والقمر.
ثم تُطفأ \omterm{def:g11:fundamental-interactions:four}{الجاذبية}: فأيّ شحنتين $+Q$ على الأرض و$-Q$
على القمر تبقيان التجاذب نفسه؟ وكم إلكترونًا هي $Q$، وكم
تزن ($m_e = \qty{9.11e-31}{kg}$)؟ علّق.
\end{exercise}

\begin{exercise}[$\star\star\star$]\label{exo:g11:fundamental-interactions:15}
في بلورة ملح، يقع أيونا Na$^{+}$ وCl$^{-}$ (وشحنتاهما $\pm e$) على بعد
$d = \qty{2.8e-10}{m}$؛ وكتلة أيون Cl$^{-}$ هي \qty{5.9e-26}{kg}.
احسب \omterm{def:g10:inertia:force}{القوة} بين الأيونين المتجاورين \omterm{def:g10:universal-gravitation:weight}{ووزن} أيون
Cl$^{-}$؛ وأعطِ النسبة. ثم اشرح كيف تفوز \omterm{def:g11:fundamental-interactions:four}{الجاذبية} على المقاييس
الكبيرة --- فلا جبل يتجاوز \qty{e4}{m} --- مع أن كل رابطة تغلبها
بخمس عشرة رتبة من القدر.
\end{exercise}

\section{مسألة: تدقيق القوى الأربع}

\begin{problem}[{مسألة نهاية الأسبوع --- لماذا لا ينهار العالم
ولا يتطاير: تدقيق التفاعلات الأربعة طابقًا طابقًا، وصولًا إلى
تعادل المادة بجزأين من $10^{18}$}]
\label{pb:g11:fundamental-interactions:1}
الذرات لا تنفجر إلى الداخل، والنوى تصمد في معظمها، والقمر لا يرتطم ولا
يفلت؛ وتدقّق هذه المسألة من يستحق الفضل. المعطيات:
$e = \qty{1.60e-19}{C}$،
$k = \qty{8.99e9}{N.m^2/C^2}$، $G = \qty{6.67e-11}{N.m^2/kg^2}$،
$m_p = \qty{1.67e-27}{kg}$، $m_e = \qty{9.11e-31}{kg}$،
$M_E = \qty{5.97e24}{kg}$، $M_M = \qty{7.35e22}{kg}$، والمسافة بين الأرض
والقمر \qty{3.84e8}{m}.

\medskip
\noindent\textbf{الجزء I --- السُّلّم.}
\begin{enumerate}
  \item أحصِ طوابق السُّلّم من \omterm{def:g11:fundamental-interactions:ladder}{الكوارك} إلى المجرّة، مع رتبة
        قدر واحدة للحجم في كلٍّ.
  \item مجسّم بمقياس: تصير \omterm{def:g11:fundamental-interactions:ladder}{النواة} بلية قطرها \qty{1}{cm}؛ فكم يبلغ عرض
        \omterm{def:g11:fundamental-interactions:ladder}{الذرة}؟
  \item وأيّ جزء من حجم \omterm{def:g11:fundamental-interactions:ladder}{الذرة} تشغله نواتها؟
  \item سُمك شعرة \qty{70}{\micro m}. فكم \omterm{def:g11:fundamental-interactions:ladder}{ذرة} تمتد عليها؟
  \item وكم \omterm{def:g10:inertia:force}{قوة} من قوى العشرة من سقف \omterm{def:g11:fundamental-interactions:ladder}{الكوارك} (\qty{e-18}{m})
        إلى مجرّة (\qty{e21}{m})؟
\end{enumerate}

\medskip
\noindent\textbf{الجزء II --- الطابق الكهربائي.}
\begin{enumerate}[resume]
  \item في \omterm{def:g11:fundamental-interactions:ladder}{ذرة} هيدروجين يقع الإلكترون على بعد $d = \qty{5.3e-11}{m}$
        من البروتون: احسب \omterm{def:g10:inertia:force}{القوة} الكهربائية بينهما.
  \item احسب \omterm{def:g10:inertia:force}{القوة} الثقالية بينهما، ونسبة
        القوتين. فأيّهما يمسك \omterm{def:g11:fundamental-interactions:ladder}{الذرة} معًا؟
  \item وماذا يتغير إذا أُطفئت \omterm{def:g11:fundamental-interactions:four}{الجاذبية} داخل الذرات؟ وإذا أُطفئت
        الكهرباء؟
  \item ولماذا تكون قوى ``التماس'' في الحياة اليومية --- كدفع الأرضية
        لقدميك، \omterm{def:g10:inertia:inventory}{والاحتكاك} --- \omterm{prop:g10:signals-and-waves:em}{كهرمغناطيسية} متنكّرة؟
  \item تحوي كل قطعة نقود من قطعتين نحو $\num{e23}$ إلكترون. انقل
        إلكترونًا واحدًا من كل \emph{مليار} بينهما، وضعهما على بعد
        \qty{1.0}{m}: احسب الشحنتين \omterm{def:g10:inertia:force}{والقوة}. فما الاستنتاج؟
  \item وأيّ تفاعل يحكم الطوابق من \qty{e-10}{m} إلى
        \qty{e4}{m}، ولماذا لا ينافسه \omterm{def:g11:fundamental-interactions:strong}{التفاعل النووي القوي}؟
\end{enumerate}

\medskip
\noindent\textbf{الجزء III --- الطابق النووي.}
\begin{enumerate}[resume]
  \item يقع بروتونان على بعد \qty{1.0e-15}{m} في \omterm{def:g11:fundamental-interactions:ladder}{نواة}: احسب
        تنافرهما الكهربائي.
  \item ولو تُرك بلا كابح، أيّ تسارع كانت هذه \omterm{def:g10:inertia:force}{القوة} ستعطيه بروتونًا؟
        قارنه بالمقدار $g$.
  \item أحصِ الخواص الثلاث التي يحتاجها \omterm{def:g11:fundamental-interactions:strong}{التفاعل النووي القوي}
        ليفسّر لماذا توجد النوى، ولماذا تُمسك النيوترونات أيضًا، ولماذا
        تكون النوى ضئيلة.
  \item في \omterm{def:g11:fundamental-interactions:ladder}{نواة} يورانيوم (قطرها \qty{1.4e-14}{m}، وفيها $92$ بروتونًا)،
        احسب التنافر بين بروتونين على طرفيها المتقابلين. فهل يقدر
        \omterm{def:g11:fundamental-interactions:strong}{التفاعل النووي القوي} على ربط ذلك الزوج مباشرة؟ واشرح لماذا
        يفوز التنافر كلما كبرت النوى، وما فائدة النيوترونات.
  \item وفي جملة واحدة، كما وُعد: ماذا يفعل \omterm{def:g11:fundamental-interactions:weak}{التفاعل النووي الضعيف}؟
\end{enumerate}

\medskip
\noindent\textbf{الجزء IV --- الطابق الفلكي، والحكم.}
\begin{enumerate}[resume]
  \item احسب \omterm{def:g10:universal-gravitation:force}{قوة الجاذبية} بين الأرض والقمر.
  \item تحوي الأرض نحو $\num{1.8e51}$ بروتون (وبعددها
        إلكترونات)، ويحوي القمر نحو $\num{2.2e49}$. فأيّ جزء $f$ من
        شحنة بروتونات كل جسم، إذا تُرك بلا إلغاء، يجعل
        \omterm{def:g10:inertia:force}{القوة} الكهربائية توافق \omterm{def:g10:inertia:force}{القوة} الثقالية؟
  \item ولماذا تحكم \omterm{def:g11:fundamental-interactions:four}{الجاذبية}، ضعيفة السؤال 7،
        علم الفلك؟ سببان.
  \item الحكم --- جملة واحدة لكل طابق (\omterm{def:g11:fundamental-interactions:ladder}{النواة}، ومن \omterm{def:g11:fundamental-interactions:ladder}{الذرة} إلى الجبل،
        والكوكب فما فوق): ما الذي يمسكه معًا؟ ثم أجِب عن العنوان:
        لماذا لا ينهار العالم ولا يتطاير؟
\end{enumerate}
\end{problem}
